\section{The game}\label{sec:game}

rstlne follows the current rules of the syndicated show. Three players play four main rounds, and the player with the most money after all 4 rounds play a bonus round alone. Appendix~A lists the full set of rules and the wheel for each round.

\subsection{Turns and money}

Each main round consists of the players solving a single puzzle. On a turn, a player can spin, buy a vowel or try to solve.
\begin{itemize}
  \item A spin that lands on a dollar value pays that value for each copy of the consonant the player called. For example, 3 copies of ``T'' at \$700 pays \$2,100.
  \item Players keep their turn when they call a letter correctly. Anything else, such as a wrong or repeated letter (except on Free Play), or an incorrect solve, passes the turn to the next player. Players can also lose their turn if the wheel lands on Lose a Turn or Bankrupt.
  \item A vowel costs \$250 regardless of if it's in the puzzle or not.
  \item Money sits in a dedicated round bank until the player solves. Only with a correct solve do players pocket that round bank, or \$1,000, whichever's higher; players who failed to solve in time lose their round bank. 
\end{itemize}

\subsection{Wheel}

The wheel has 24 equally sized wedges that players spin. On all 4 rounds, there are 2 fully sized Bankrupts and 1 Lose a Turn. As the name implies, Bankrupts wipe the round bank (only) and any cards held and passes the turn, whereas Lose a Turn only passes the turn. Additionally, each round consists of some special wedges, listed below, as well as a singular top dollar wedge (\$2,500 in round 1, \$3,500 in round 2 and 3, and \$5,000 in round 4). Every other wedge is a normal dollar value from \$500 to \$900. 

\begin{itemize}
  \item The Wild Card (rounds 1 to 3) grant players an extra consonant to call, usable either during the main or bonus round.
  \item The Million wedge (rounds 1 to 3) can turn the max possible bonus prize to \$1,000,000. It only takes up 1/3 of the wedge it's on; the surrounding 2/3's are Bankrupts.
  \item The Mystery wedges (round 2) have a face value of \$1,000. Players can keep it or flip it, which can reveal either \$10,000 or Bankrupt. There are 2 copies of this wedge.
  \item The Express (round 3) also have a face value of \$1,000. Players can opt to ``board'' the Express, meaning that they are obligated to continuously call correct letters until they solve successfully. Consonants pay \$1,000 per copy and a vowel buy is \$250 as usual. Any outcome besides a successful solve (e.g. failed solve, missed letter, etc.) terminates their run and bankrupts the player. They can also pass.
  \item The Free Play (all 4 rounds) allow players to call any letter of their choosing and keeps their turn even if their letter misses. Consonants count for \$500 apiece and vowels are free. Although this wedge was removed since 2021's iteration of the show, we kept this wedge to see how players make choices under this state.
\end{itemize}

\subsection{Bonus round}

Bonus rounds, like the syndicated show, open with R, S, T, L, N, and E showing. The player picks 3 more consonants and 1 more vowel (4 consonants with a Wild Card). They then get one chance to solve the puzzle, with the right answer winning an envelope worth between \$40,000 to \$100,000 (also determined by a wheel spin with the value initially hidden), or \$1,000,000 replacing the \$100,000 envelope with the Million wedge on hand.

\subsection{Deviations from the show}

\begin{itemize}
  \item There are no toss ups. Toss ups work by progressively revealing letters at random on a puzzle, with each of the 3 players giving one chance to solve. The first to get the puzzle right wins the preset dollar value of that toss up. Toss ups are innately speed rounds and we are unable to implement speed in this project. We considered polling every player on when to solve for each letter reveal, but that removes the speed aspect and favors non-human players like our bot and the LLMs.
  \item There are no clocks. Each player is allowed to take however long they need to make their move.
  \item Seats 1, 2 and 3 always start rounds 1, 2 and 3, with the winner of round 3 starting round 4. On the show, toss ups decide who starts rounds 1 and 4.
  \item Free Play stays on the wheel despite the show retiring it in 2021.
  \item On the bonus round, players on the show get 10 seconds to guess the puzzle, but they can guess as many times as they want (and some of them think aloud). Players in this project get only one guess but can spend however long they want to submit it. Additionally, contestants pick the bonus puzzle category in the show, whereas this is randomly picked in this project.
  \item Additionally, the canonical show contains sponsored prizes such as gift tags, prize wedges (earning players a trip), a ``prize puzzle'' for round 3, and potential cars on some envelopes in the bonus round. They are all left out in this project so the players play only for money.
\end{itemize}

\subsection{The game engine}

The engine simply faithfully implements the design of this game in accordance to the rules and deviations listed above. For the sake of reproducibility, the PRNG (pseudorandom number generator) is pre-seeded such that reruns always yield the same sequence of wheel outcomes and puzzles. To support LLMs, we implement moves as tool calls. Incorrect tool calls are rejected with reasons given to them, and a third consecutive incorrect tool call passes the turn. Players only see what contestants see on the canonical show to avoid cheating (so no advance knowledge of the puzzles or peeks at Mystery wedges or bonus envelopes). Section~\ref{sec:can-they-play} lists them in detail.

